% Источники: exam/materials/преподаватель/2020/23-03-2020-MODEL L5(1).doc; exam/materials/моделирование.pdf, раздел 11; exam/materials/прошлогодние ответы/Modelirovanie_1_-3.pdf, вопрос 16; https://github.com/HumsterProgrammer/iu9-notes/blob/main/8sem/Моделирование/Моделирование_лекция-09_17-03-26.md.
\section{Стохастическое моделирование. Доверительное оценивание параметров модели.}

\textbf{Доверительный интервал} для параметра $\theta$ — интервал $[\hat{\theta}_\text{н},\, \hat{\theta}_\text{в}]$, построенный по выборке, такой что
\[
  P\bigl(\hat{\theta}_\text{н} \leqslant \theta \leqslant \hat{\theta}_\text{в}\bigr) = \gamma,
\]
где $\gamma$ — \textbf{доверительная вероятность} (уровень доверия), $1-\gamma$ — уровень значимости.

\subsection*{Построение доверительного интервала}

Метод опирается на \textbf{центральную статистику} $T = T(x_1,\ldots,x_n;\,\theta)$ — функцию выборки и параметра, распределение которой не зависит от $\theta$.

Для $T$ находят квантили $t_{\varepsilon_1}$ и $t_{1-\varepsilon_2}$ такие, что
\[
  P\bigl(t_{\varepsilon_1} \leqslant T \leqslant t_{1-\varepsilon_2}\bigr) = \gamma,
\]
где $\varepsilon_1 + \varepsilon_2 = 1 - \gamma$.

При \textbf{симметричном} интервале:
\[
  \varepsilon_1 = \varepsilon_2 = \frac{1-\gamma}{2}.
\]

Из неравенства $t_{\varepsilon_1} \leqslant T \leqslant t_{1-\varepsilon_2}$ выражают $\theta$ и получают $[\hat{\theta}_\text{н},\, \hat{\theta}_\text{в}]$.

\subsection*{Доверительный интервал для среднего (нормальное распределение)}

Если $\sigma^2$ известна:
\[
  \bar{x} - z_{1-\varepsilon}\,\frac{\sigma}{\sqrt{n}} \leqslant \mu \leqslant \bar{x} + z_{1-\varepsilon}\,\frac{\sigma}{\sqrt{n}},
\]
где $z_{1-\varepsilon}$ — квантиль стандартного нормального распределения.

Если $\sigma^2$ неизвестна (заменяется $s$): используют распределение Стьюдента $t_{n-1}$:
\[
  \bar{x} \pm t_{1-\varepsilon,n-1}\,\frac{s}{\sqrt{n}}.
\]

\subsection*{Влияние объёма выборки}

Ширина интервала $\hat{\theta}_\text{в} - \hat{\theta}_\text{н} \propto 1/\sqrt{n}$: для уменьшения ширины вдвое необходимо увеличить число испытаний вчетверо.

\clearpage
